\documentclass[11pt]{article}
\usepackage{tcolorbox}
\usepackage{textcomp}
\usepackage[os=win]{menukeys}
\usepackage{pdflscape}
\usepackage{rotating}
\usepackage{listings}
\usepackage{titling}
\usepackage{enumitem}
\setlist{nosep}


\newcommand{\numpy}{{\tt numpy}}    % tt font for numpy

\topmargin -.5in
\textheight 9in
\oddsidemargin -.25in
\evensidemargin -.25in
\textwidth 7in

\pretitle{%
  \begin{center}
  \LARGE
  \includegraphics[width=15cm]{../images/logo.png}\\[\bigskipamount]
}
\posttitle{\end{center}}

\begin{document}

% ========== Edit your name here
\author{Dr. Neil Eliot / Dr. Alun Moon}
\title{KF5004\\------\\Business case, scenario and the exam\\------\\The Implementation of a \texttt{DNS} infrastructure (\texttt{BIND9}) 
and Web Server Farm facility for a large scale organisation.
(\texttt{Apache, PHP, NFS, and MySQL})\\------}
\date{September 2019}
\maketitle

\newpage
\tableofcontents
\newpage

\section{Introduction}
The assessment of \texttt{KF5004} is an exam and accounts for 100\% of the module mark. It is based on a business case that you will be working on 
over the remainder of the semester.\\

\noindent The business scenario is supplied in a separate document. 

\section{Schedule}
By week 12 you should be in a position to demonstrate a full implementation. Over weeks 10-12 you will be developing a solution for the company. 
A sample schedule of work is shown in table 1.

\begin{table}[ht]\label{schedule}
    \begin{tabular}{|p{1cm}|p{5cm}|p{5cm}|p{5cm}|} 
      \hline
      \multicolumn{1}{|c|}{Week} & \multicolumn{1}{|c|}{Lab Activity 1} & \multicolumn{1}{|c|}{Lab Activity 2}& \multicolumn{1}{|c|}{Additional} \\ 
      \hline
      10a & Lecture/Workshop& Briefing Out & Create Multiple staticlly addressesd servers\\
      \hline
      10b & Assignment briefing.& Attempt your own design.& Analyse provided design and build base servers (basic service installations)\\
      \hline
      11a & Setup \texttt{DNS} architecture.& Basic \texttt{Apache} installs. & Test load balanced deployment and upload sample websites\\
      \hline
      11b & Database setup, config and test connectivity. (Intranet and Internet)& \texttt{NFS} deployment.& Connectivity testing of Database and \texttt{NFS}\\
      \hline
      12a & Feedback & Demos & Updates\\
      \hline
      12b & Demos & Demos & Revision\\
      \hline
    \end{tabular}
    \caption{Activity Schedule}
\end{table}
\newpage 

\section{Workload}
A breakdown of the ``effort'' required for each section is detailed in Table 2.

\begin{table}[ht]\label{workload}
    \begin{tabular}{|p{1cm}|p{14cm}|p{1cm}|} 
      \hline
      \multicolumn{1}{|c|}{Task} & \multicolumn{1}{|c|}{Description} & \multicolumn{1}{|c|}{Effort (\%)}\\ 
      \hline
      1 & \texttt{DNS} Core Configuration 
            \begin{itemize}
                \item SOA Record setup
                \item Other Record setup
                \item DNS Core Operations 
            \end{itemize}
        & 20\%\\
      \hline
      2 & \texttt{Apache} (Intranet) Configuration
            \begin{itemize}
                \item Intranet site including scripting and database connectivity
                \item Personal web storage
                    \begin{itemize}
                        \item \texttt{FTP} facilities
                        \item Static sites
                        \item Dynamic site
                    \end{itemize}
            \end{itemize} 
        & 20\%\\
      \hline
      3 & Database (\texttt{Intranet}) Configuration
            \begin{itemize}
              \item Intranet site
              \item Management console
              \item Localized security    
            \end{itemize}
      & 10\%\\
      \hline
      4 & Apache (\texttt{Internet}) Configuration including Load Balancing
            \begin{itemize}
              \item Individual server access
              \item Load-balanced server access
              \item Internet site including scripting and database connectivity
              \item Localised scripting
            \end{itemize} 
      & 20\%\\
      \hline
      5 & Content Storage Configuration (for web server farm)
          \begin{itemize}
            \item \texttt{NFS} configuration
            \item \texttt{rDNS} Security
          \end{itemize} 
      & 20\%\\
      \hline
      6 & Database (\texttt{Internet}) Configuration (for web server farm)
          \begin{itemize}
            \item Management Console (\texttt{PHPMyAdmin})
            \item Database creation with user accounts
            \item Remote access security
          \end{itemize}
      & 10\%\\
      \hline
    \end{tabular}
    \caption{Activity Schedule}
\end{table}

\section{Examination}

\noindent The examination is be based on the business scenario you have been provided with but will use a different set of \texttt{IP} 
addresses to your block allocation. At the examination you will be provided with both the scenario and the possible configuration.\\

\noindent There will be 15 mins reading time at the beginning of the examination for you to familiarise yourself with the configurations. 
Following the reading time, you will have a multi-choice examination based on both the teaching in the module and the scenario.\\

\begin{tcolorbox}[title={\textbf{NOTE:}}]
  \noindent \textbf{THE SAMPLE CONFIGURATION AND THE EXAM QUESTIONS MUST NOT BE REMOVED FROM THE EXAM ROOM.}
\end{tcolorbox}

\section{Feedback}

Formative (non assessed) feedback will be provided on your progress throughout the workshops as you develop your solution and at a final demonstration if you choice to participate. 
The examination results will be released via blackboard 3 weeks after the examination.\\

\noindent \texttt{END}
\end{document}
